\Needspace{7\baselineskip}
\section{Observer and Readout Architecture}
\label{sec:08}
Section 7 closed the internal junction account. It proved the internal update-closure identity, fixed the closure-side status, and stopped before observer-system readout. Section 8 begins on the other side of that boundary. Its task is not to import measured physics. Its task is to define what a readout locus is, how completed internal records become available to that locus, how those registrations function as a clock, and how completed closure extent becomes available as a time readout without changing the internal book.

The dependency direction is one-way. The readout layer consumes the completed internal record; it does not ground, reopen, or edit it. The word "observer" is therefore not a psychological word here. It names a relational role: the standpoint through which already resolved distinctions are available for reading.

\Needspace{5\baselineskip}
\subsection{The Readout Locus}
The first object is the readout locus \(L\). A locus is a relational standpoint at an internal resolution scale. Its content is the set of already resolved completed updates distinguishable at that scale:

\begin{equation*}
\begin{adjustbox}{max width=\linewidth,center}
$\displaystyle
\mathcal{D}(L)=\{u:\operatorname{reg}(L,u)\}.
$
\end{adjustbox}
\end{equation*}
% LAYOUT_ONLY_GUARD:REGISTERED_ELEMENT_PARAGRAPH
\Needspace{10\baselineskip}
The elements \(u\) are completed internal updates. They are registrable only after the internal resolution condition has been met; below that floor there is no registered object, not a stored zero value. Each registered element still carries its orientation. The locus does not create those distinctions, does not resolve them, and does not duplicate the aperture that resolved them. It only makes already resolved content available on the readout side.

This first set has no order. It is a membership surface: what is available at \(L\), not what came first, what follows, how much time has elapsed, or what physical rate is being measured. Those are later structures.

\Needspace{5\baselineskip}
\subsection{Registration}
Registration is the relation that populates the readout set. For a locus \(L\) and an already resolved update \(u\), the statement

\begin{equation*}
\begin{adjustbox}{max width=\linewidth,center}
$\displaystyle
\operatorname{reg}(L,u)
$
\end{adjustbox}
\end{equation*}
means that \(u\) is available at \(L\). Equivalently, \(u\in\mathcal{D}(L)\). This is a yes-or-no membership relation. It carries the orientation of \(u\), but it does not impose a sequence, compare positions, assign a duration, create a history of registrations, or turn the set into an ordered measure.

The distinction matters because a registration relation can be mistaken for measurement. It is not measurement in the later physical sense. It is availability of already resolved relational content at a locus. The internal structure has already done the resolving; registration reads it.

\Needspace{5\baselineskip}
\subsection{The Reference Role}
A distinguishable reading also needs a reference. For an existent or system \(X\) at a level \(\ell\), the observer/reference role may be written as

\begin{equation*}
\begin{adjustbox}{max width=\linewidth,center}
$\displaystyle
\operatorname{Obs}(X,\ell).
$
\end{adjustbox}
\end{equation*}
The role holds when \(X\) occupies a readout standpoint, supplies the relational reference relative to which registered content is distinguishable, and thereby makes a distinguishable readout available at that level. The reference is not a background frame and not a privileged object. It is a relation. A registered state is distinguishable from the standpoint that reads it.

This is the sense in which observerhood induces readout asymmetry. The observer does not create distinguishability. The distinguishability is primitive and the relevant completion has already been resolved. By occupying the reference role, \(X\) selects which oriented side is read as reference. The resulting asymmetry is an orientation selection, not a signed amount and not an act of creation.

The role is independent of a fixed cell count. One cell may occupy the standpoint. Without a second registered state there is no realized distinction to read against it; with more registered states there is more distinguishable content. The role supplies the reference condition, not a cell-number theorem.

\Needspace{5\baselineskip}
\subsection{Clockhood}
The unordered set \(\mathcal{D}(L)\) is still not a clock. Clockhood is a role in which registered content is read as an ordered measure. It requires three ingredients.

First, there must be a succession, inherited from the completed-closure count. Second, there must be a relational comparison standard against which that succession is read. Third, the locus reads the succession against the standard as a relative, pre-metric ordered measure.

The order therefore lives in the reading, not in the set. The set remains unordered; registration remains membership. A wristwatch, recurrence, heartbeat, or astronomical cycle is clock-like only when it occupies this role. Nothing is inherently a clock.

The comparison standard is relational. It is not a background grid, a privileged frame, a fixed external coordinate system, or an absolute clock. A recurrence may serve, but a second recurrence is not required; a single relational comparison baseline can supply the standard. Clockhood gives readable order. It still gives no physical unit, no calibrated time, no measured rate, and no prediction.

\Needspace{5\baselineskip}
\subsection{The Time Readout}
Through clockhood, the locus reads the inherited succession of registered completions against its relational comparison standard. Clockhood creates neither the completions nor their geometry. The time readout adds the per-closure extent. Each completed closure carries the dimensionless base extent

\begin{equation*}
\begin{adjustbox}{max width=\linewidth,center}
$\displaystyle
t_{\mathrm{base}}=2\pi .
$
\end{adjustbox}
\end{equation*}
Thus the time readout \(\mathrm{t}(L)\) is the locus reading its registered completed closures as carrying that per-closure extent. In closure units, it is the completed-closure count together with \(2\pi\) of extent per closure:

\begin{equation*}
\begin{adjustbox}{max width=\linewidth,center}
$\displaystyle
\mathrm{t}(L)=I\,t_{\mathrm{base}}.
$
\end{adjustbox}
\end{equation*}
This is why the time readout is not merely clockhood renamed. Clockhood supplies readable order; \(\mathrm{t}(L)\) binds the \(2\pi\) extent to each completed closure. The result is dimensionless, metric-free, and not yet calibrated into seconds or any other observer-scale unit. The later calibration relation is only a pointer at this stage; it is not folded into the readout.

\Needspace{5\baselineskip}
\subsection{Book-Energy Record}
Separately from the time readout, the established bookkeeping coefficient and channel-route count give the internal book-energy record

\begin{equation*}
\begin{adjustbox}{max width=\linewidth,center}
$\displaystyle
\mathrm{Cap}=c^{\otimes 2}=9,
\qquad
E=m\cdot \mathrm{Cap}=\frac{9}{2\pi}.
$
\end{adjustbox}
\end{equation*}
Here \(m\) is the active bookkeeping coefficient used in this expression. The expression is the channel-distribution image of that coefficient across the forced channel-route count. The symbol \(c\) in \(\mathrm{Cap}=c^{\otimes 2}\) is a channel-count factor, not the measured speed of light. The value \(9\) is the channel-route count, not a physical velocity squared. No physical rate or calibrated energy law is inferred from this internal record.

\Needspace{5\baselineskip}
\subsection{Structural Completion and Prediction Boundary}
At this point the observer/readout stack has an internal structural account of its source record, readout locus, registration, reference role, clock target, closure extent, and dependency direction. External calibration and comparison remain deferred.

That status does not unlock measured relativity, Lorentz recovery, time dilation, measured rates, empirical prediction, or gravity. The section therefore records internal structural completeness while refusing the measured-physics unlock.

The prediction layer is opened only procedurally. Later work must first derive the theory's own comparison-law analogues and only then compare externally. A measured prediction remains inadmissible until the relevant comparison law and measured bridge conditions have been supplied. The internal update-closure identity from Section 7 is available as internal theorem ground, but it cannot be converted into physical readout, empirical comparison, or prediction by this section.

\Needspace{5\baselineskip}
\subsection{Boundary of the Section}
Section 8 constructs the observer/readout architecture. It defines the readout locus, registration relation, observer/reference role, clockhood reading of inherited succession, closure-unit time readout, and the separately established internal book-energy record. It opens external calibration and prediction only as ordered future obligations.

It does not import measured time, assign SI units as fundamental, derive a measured mass-energy law, derive Lorentz recovery, derive gravity, produce a pre-comparison output, or claim empirical success.

This layer can also lose, and says how. Exhibit a readout locus that creates rather than consumes distinguishability, and the locus definition fails. Exhibit registration that orders the set, and the registration definition fails. Exhibit a clock without a relational comparison standard, and clockhood fails. Exhibit a time readout without \(2\pi\) per closure, and the time readout fails. Treat the internal book-energy record as a calibrated physical law, and its scope fence fails. Exhibit a measured prediction before the internally derived comparison law and measured comparison conditions exist, and the prediction boundary fails.

